\subsection{Non-Abelian Discrete Flavor Symmetries from Magnetized/Intersecting Brane Models --- arXiv:0904.2631}

\textbf{Authors:} Hiroyuki Abe, Kang-Sin Choi, Tatsuo Kobayashi, Hiroshi Ohki

\paragraph{主な主張}
磁化・交差D-brane模型の零モード結合から$D_4$、$\Delta(27)$、$\Delta(54)$などの非可換離散対称性と物質の表現を導き、Wilson lineやオービフォールド射影による変化を示す~\cite{Abe:2009vi}。

\paragraph{新規性}
波動関数の重なり積分の選択則と零モード間の変換から、非可換フレーバー対称性の具体的起源を明らかにする。

\paragraph{位置付け}
磁化余剰次元の零モード構造をフレーバー対称性へ結び付ける初期の基礎文献。

\paragraph{コメント（読書メモ）}
本人の読書メモ：arXiv:1404.0137~\cite{Abe:2014nla}で、離散対称性を得る方法として引用・レビューされていた。

Wilson lineがない場合について、磁束の整数の最大公約数を$g$とすると、トーラス上の重なり積分からYukawa結合に$\mathbb{Z}_g$選択則が現れ、零モードの量子数がmod $g$で制御されると整理した。

Wilson lineを含む場合の例として、$g=2$で$D_4$、$g=3$で$\Delta(27)$が現れる点を記録した。
